% !TeX program = lualatex
% Source-PDF: 比赛 CONTEST/2021“MINIEYE杯”中国大学生算法设计超级联赛/2021“MINIEYE杯”中国大学生算法设计超级联赛（7）/solution.pdf
\documentclass[11pt]{article}
\usepackage{oipl-vn}

\title{Lời giải trận thứ bảy HDU Multi-University 2021}
\author{Ban to chuc}
\date{}

\begin{document}
\maketitle

\origtitle{2021 ``MINIEYE Cup'' China Collegiate Programming Contest Training, Contest 7 Solutions}
\sourcepath{contest/hdu-2021-minieye-07-solutions.pdf}

\section{Fall with Fake Problem}

Xét việc nhị phân vị trí của chữ cái đầu tiên trong chuỗi $T$ khác với chuỗi
$S$.

Trong quá trình \texttt{check}, ta muốn dựng một chuỗi có thứ tự từ điển lớn
nhất có thể. Khi đó, từ vị trí đang \texttt{check} trở đi, chuỗi $T$ chắc chắn
là một dãy giảm. Xét vị trí thật sự đầu tiên khác với chuỗi gốc\footnote{Ví
dụ, với \texttt{zzzyy}, khi \texttt{check(2)} thì vị trí đầu tiên khác chuỗi
gốc ít nhất là vị trí thứ $4$.}, số vị trí có thể nhiều nhất là $25$
\footnote{Từ vị trí nhị phân hiện tại, chuỗi $T$ chắc chắn là một dãy giảm;
nếu không, ta có thể đưa chữ cái phía sau lên trước để thứ tự từ điển lớn hơn.
Vì vậy dãy chỉ có thể giảm nhiều nhất $25$ lần, và vị trí đầu tiên khác chuỗi
gốc chỉ có thể là lần xuất hiện đầu của mỗi loại chữ cái trong dãy giảm.}.
Giả sử chữ cái đầu tiên trong $T$ khác với $S$ là $c$. Khi đó vị trí của mọi
chữ cái lớn hơn $c$ đã được xác định\footnote{Nếu sau vị trí này vẫn còn chữ
cái lớn hơn $c$, ta có thể chuyển chữ cái đó lên trước để được chuỗi có thứ tự
từ điển lớn hơn.}. Ta có thể liệt kê số lần xuất hiện của $c$\footnote{Số lần
xuất hiện là một ước của $k$, nên không nhiều; trong độ phức tạp ký hiệu là
$\sqrt{k}$.}; các vị trí còn lại cần được lấp bằng các chữ cái nhỏ hơn $c$.
Phần này có thể giải bằng ba lô vét cạn, dùng \texttt{bitset} để tăng tốc.
Độ phức tạp của \texttt{check} đạt
\[
  O\!\left(25\sqrt{k}+25\sqrt{k}\frac{n}{64}\right).
\]
Thực tế quá trình liệt kê $c$ không đầy, vì số lần xuất hiện của các chữ cái
lớn hơn $c$ thường không phải là ước của $k$.

Khi dựng đáp án, phần đầu của chuỗi $T$ đã được xác định bởi các bước nhị phân
trước đó; phần sau cũng dùng ba lô tăng tốc bằng \texttt{bitset}, rồi xuất một
nghiệm nhỏ nhất theo quá trình DP. Cách làm chuẩn liệt kê chữ cái cụ thể đóng
vai trò chữ cái đầu tiên khác nhau, sau đó chạy DP giống trong \texttt{check}
và xuất nghiệm theo quá trình DP. Độ phức tạp là
\[
  O\!\left(25\sqrt{k}\frac{n}{64}\right).
\]
Độ phức tạp cho một bộ kiểm thử là
\[
  O\!\left(25\sqrt{k}\frac{n}{64}\log n\right).
\]

\section{Fall with Soldiers}

Trong bản lời giải ngắn, chỉ cần đọc các phần chữ đậm.

Trước hết chú ý tính chất then chốt của đề: \term{độ dài chuỗi là số lẻ}. Mỗi
lần thao tác ta xóa hai ký tự hai bên của mỗi ký tự $1$, nên tính chẵn lẻ của
độ dài chuỗi không đổi. Dễ thấy nếu có một phương án thao tác làm cho cuối cùng
chuỗi còn lại toàn là $1$, thì ta có thể tiếp tục xóa chuỗi còn lại cho đến khi
chỉ còn một ký tự $1$. Ngược lại, nếu có phương án xóa chuỗi đến khi chỉ còn
một ký tự $1$, phương án đó chắc chắn thỏa điều kiện đề bài. Vì vậy điều kiện
đề bài cần tìm tương đương với điều kiện có thể xóa chuỗi đến khi chỉ còn một
ký tự $1$.

Bây giờ xét khi nào một chuỗi có thể bị xóa đến chỉ còn một ký tự $1$. Nếu vị
trí giữa chuỗi là $1$, ta có thể liên tục thao tác trên ký tự $1$ này để đạt
mục tiêu. Nếu vị trí giữa không phải là $1$, ta xét điều kiện để đoạn $0$ ở
giữa có thể bị xóa. Mỗi ký tự $1$ trong chuỗi có số lần thao tác tối đa bằng
khoảng cách từ nó đến biên gần nhất; ví dụ \texttt{0010000} nhiều nhất chỉ thao
tác được hai lần. Một kết luận rất hiển nhiên là ta nên chọn ký tự $1$ càng gần
tâm càng tốt. Lấy hai ký tự $1$ gần tâm nhất ở hai phía trái và phải, rồi xét
các đoạn mà thao tác trên từng ký tự có thể xóa. Khi hai đoạn xóa phủ toàn bộ
chuỗi, chắc chắn tồn tại một phương án thao tác đưa toàn bộ chuỗi về chỉ còn
một ký tự $1$\footnote{Ví dụ \texttt{0010010}: phạm vi bị xóa bởi hai ký tự
$1$ trái và phải lần lượt là \texttt{DD1DD10} và \texttt{0010D1D}. Hai phạm vi
này chồng lên nhau, nên ta có thể hoàn thành mục tiêu. Một cách rõ ràng là
liên tục thao tác trên một trong hai ký tự $1$ cho đến khi ký tự $1$ còn lại
nằm ở chính giữa chuỗi, rồi thao tác trên nó. Mỗi lần thao tác trên một ký tự
$1$ làm vị trí giữa dịch về phía ký tự $1$ kia một ô; hai phạm vi xóa phủ nhau
nghĩa là ký tự $1$ kia luôn có thể đi tới vị trí giữa, nên phương án khả thi.}.
Nếu hai đoạn không phủ hết chuỗi, chắc chắn còn một ký tự $0$ nằm giữa hai ký
tự $1$ và không thể bị xóa bằng bất kỳ cách nào\footnote{Ví dụ \texttt{0010001}:
phạm vi xóa của hai ký tự $1$ là \texttt{DD1DD01} và rỗng; có vị trí
\texttt{00100X1} ở giữa không được ký tự $1$ nào phủ. Nếu tồn tại ký tự $1$
khác, phạm vi của chúng cũng không thể lớn hơn hai ký tự $1$ gần giữa nhất, nên
vị trí đó không thể bị xóa.}. Điều kiện cuối cùng là: \term{hai ký tự $1$ gần
nhất ở hai phía của đường giữa có khoảng cách không vượt quá $(n-3)/2$}.

Sau khi có điều kiện, ta có thể đếm. Gọi tổng số vị trí cần điền là
\textit{tot}. Chia trường hợp:
\begin{itemize}
  \item Nếu chính giữa là $1$, mọi vị trí còn lại điền tùy ý, đáp án là
  $2^{\textit{tot}}$.
  \item Nếu chính giữa không phải là $1$, ta chia khoảng theo đường giữa thành
  hai nửa trái và phải. Gọi $\textit{cntl}_i$ là số dấu hỏi bên trái vị trí
  thứ $i$ của nửa trái, $\textit{cntr}_i$ là số dấu hỏi trong nửa phải có
  khoảng cách đến vị trí đó không vượt quá $(n-3)/2$, $\textit{locl}$ là vị trí
  ký tự $1$ ngoài cùng bên phải trong nửa trái, $\textit{locr}$ là vị trí ký tự
  $1$ ngoài cùng bên trái trong nửa phải, và $\textit{totr}$ là tổng số dấu hỏi
  trong nửa phải. Xét liệt kê vị trí ký tự $1$ ngoài cùng bên phải của nửa trái:
  \begin{itemize}
    \item với $i<\textit{locl}$, không có cách điền dấu hỏi thành $1$ nào thỏa
    điều kiện, nên không có đóng góp;
    \item với $i\ge \textit{locl}$, để thỏa điều kiện thì vị trí này phải là
    $1$, mọi vị trí bên phải nó trong nửa trái phải điền $0$\footnote{Trường
    hợp này gồm cả khi điểm hiện tại là dấu hỏi và khi điểm hiện tại chính là
    $\textit{locl}$, nhưng cách tính giống nhau.}, và trong nửa phải phải có ít
    nhất một ký tự $1$ thỏa giới hạn khoảng cách. Xét quan hệ giữa
    $\textit{locr}$ và vị trí đó, ta nhận được hai dạng đóng góp
    \[
      2^{\textit{cntl}_i+\textit{totr}-\textit{cntr}_i}
      \left(2^{\textit{cntr}_i}-1\right)
      =
      2^{\textit{cntl}_i+\textit{totr}}
      -
      2^{\textit{cntl}_i+\textit{totr}-\textit{cntr}_i},
    \]
    hoặc
    \[
      2^{\textit{cntl}_i+\textit{totr}}.
    \]
  \end{itemize}
\end{itemize}

Dễ thấy ta cần duy trì, với mỗi vị trí dấu hỏi trong nửa trái, hai giá trị
$2^{\textit{cntl}_i+\textit{totr}}$ và
$2^{\textit{cntl}_i+\textit{totr}-\textit{cntr}_i}$. Đối với đáp án tại
$\textit{locl}$, còn cần duy trì $\textit{cntl}_i$ và $\textit{cntr}_i$. Khi
thay đổi ký tự ở một vị trí, thông tin trên một đoạn sẽ bị ảnh hưởng bằng các
phép nhân/chia $2$ hoặc cộng/trừ $1$. Ta có thể dùng cây đoạn để duy trì những
thông tin này và khi truy vấn thì phân loại như trên để tính đáp án.

\section{Fall with Trees}

Đặt độ rộng tầng thứ $2$ là $w_2=d$. Khi đó độ rộng tầng thứ $k$ là
\[
  w_k=\sum_{i=2}^{k} d\cdot 2^{2-i}
      = d\left(2-2^{2-k}\right),
\]
và diện tích giữa tầng $k$ và tầng $k+1$ là
\[
  S_k=h\frac{w_k+w_{k+1}}{2}.
\]
Vì vậy ta cần tính
\[
  S=\sum_{i=1}^{k-1}S_i
   =\frac{hd}{2}\sum_{i=1}^{k-1}\left(4-3\cdot 2^{1-i}\right).
\]
Phần này có thể cộng trực tiếp đến khi đạt độ chính xác yêu cầu, hoặc dùng
công thức tổng cấp số nhân.

\section{Link with Balls}

Gộp hộp có thể lấy $0\sim k-1$ quả bóng với hộp chỉ có thể lấy bội số của $k$
quả bóng thành một hộp có thể lấy số bóng tùy ý. Khi đó ta có $n$ hộp có thể
lấy tùy ý và một hộp có thể lấy $0\sim n$ quả bóng. Liệt kê số bóng lấy trong
hộp $0\sim n$; số cách chọn các quả còn lại được tính bằng phương pháp chia
thanh. Do đó đáp án là
\[
  \sum_{i=0}^{n} \binom{m-i+n-1}{n-1}
  =
  \binom{m+n}{n}-\binom{m-1}{n}.
\]

\section{Link with EQ}

Gọi $f(x)$ là kỳ vọng số người sẽ ngồi trong một đoạn trống liên tiếp độ dài
$x$ bị chặn ở cả hai đầu\footnote{Bị chặn ở cả hai đầu nghĩa là bên trái vị
trí trống ngoài cùng bên trái và bên phải vị trí trống ngoài cùng bên phải đều
đã có người.}. Khi đó:
\[
  f(x)=
  f\!\left(\left\lfloor\frac{x-1}{2}\right\rfloor\right)
  +
  f\!\left((x-1)-\left\lfloor\frac{x-1}{2}\right\rfloor\right)
  +1,
\]
\[
  f(1)=f(2)=0.
\]

Tương tự, đặt $g(x)$ là kỳ vọng số người sẽ ngồi trong một đoạn trống liên
tiếp nửa mở độ dài $x$\footnote{Nửa mở nghĩa là một phía là đầu bàn, phía còn
lại đã có người.}. Khi đó
\[
  g(x)=
  \begin{cases}
    0, & x\le 1,\\
    f(x-1)+1, & x\ge 2.
  \end{cases}
\]

Cuối cùng, gọi $h(x)$ là kỳ vọng số người sẽ ngồi ở một bàn độ dài $x$. Ta có
\[
  h(x)=\frac{1}{x}\sum_{i=1}^{x}\bigl(g(i-1)+g(x-i)+1\bigr),
\]
nên
\[
  h(x)=1+\frac{2}{x}\sum_{i=1}^{x-1}g(i).
\]
Vì vậy chỉ cần tiền xử lý lần lượt $f(x)$, $g(x)$, tổng tiền tố của $g(x)$ và
$h(x)$; khi truy vấn thì xuất trực tiếp đáp án.

\section{Link with Grenade}

Chọn ngẫu nhiên điểm trên một mặt cầu và tích phân hai lần theo góc về bản chất
là giống nhau, vì hai cách này có cùng mật độ diện tích trên đơn vị mặt cầu.
Hướng của lựu đạn trên mặt phẳng ngang không ảnh hưởng đến việc nó có nổ trúng
người hay không, nên ta chỉ cần xét mặt phẳng thẳng đứng.

Xét đổi hệ quy chiếu: cộng cho cả người và lựu đạn một gia tốc trọng trường
ngược chiều. Khi đó lựu đạn chuyển động thẳng đều, còn người chuyển động thẳng
biến đổi đều. Sau $t$ giây, vị trí của người là cố định, còn vị trí của lựu đạn
là một đường tròn; miền mà lựu đạn có thể nổ trúng người cũng là một đường
tròn. Dựng tam giác như trong hình của lời giải gốc, dùng định lý cos để tìm
góc nâng $\theta$ mà lựu đạn có thể nổ trúng người.

Theo công thức diện tích chỏm cầu, diện tích phần có thể nổ trúng người là
\[
  2\pi (v_0t)^2(1-\cos\theta),
\]
còn diện tích mặt cầu là $4\pi(v_0t)^2$. Vì vậy đáp án là
\[
  \frac{1-\cos\theta}{2}.
\]
Thay biểu thức $\cos\theta$ nhận được từ định lý cos vào là xong; chú ý xử lý
riêng hai trường hợp không thể nổ trúng và chắc chắn nổ trúng.

\section{Link with Limit}

Dựng đồ thị, nối cạnh từ đỉnh $x$ đến đỉnh $f(x)$. Khi đó quá trình lặp hàm
$f$ về bản chất là quá trình đi trên đồ thị. Bài gốc thực ra hỏi liệu trung
bình trọng số đỉnh của mọi chu trình có giống nhau hay không. Chỉ cần dùng bất
kỳ phương pháp nào để tìm tất cả chu trình và tính trung bình của chúng.

\section{Smzzl with Greedy Snake}

Trước hết, số lượng lệnh \texttt{f} hiển nhiên là khoảng cách Manhattan giữa
hai điểm kề nhau; ta chỉ cần giảm số lần rẽ. Đề bài nói hai điểm kề nhau không
nằm trên cùng hàng hay cùng cột, nghĩa là ta cần đi theo hai hướng kề nhau. Nếu
hướng hiện tại đã là một hướng cần đi, cứ đi thẳng. Nếu không, chỉ cần xoay một
lần là đến được một hướng cần đi, sau đó mô phỏng là được.

\section{Smzzl with Safe Zone}

Đề bài yêu cầu tìm một điểm trên bao lồi ngoài sao cho khoảng cách gần nhất từ
nó đến bao lồi trong là xa nhất. Có thể chứng minh rằng chọn điểm trên bao lồi
ngoài tại một đỉnh không bao giờ tệ hơn.

Do đó xét liệt kê từng điểm trên bao lồi ngoài, rồi vét cạn từng cạnh của bao
lồi trong để tính khoảng cách nhỏ nhất và lấy giá trị lớn nhất. Cách này có độ
phức tạp $O(n^2)$ cho một bộ dữ liệu.

Ta có thể dùng quay thước kẹp để tối ưu quá trình trên. Khi liệt kê các điểm
trên bao lồi ngoài theo chiều ngược kim đồng hồ, điểm gần nhất tương ứng trên
bao lồi trong cũng quay đơn điệu theo chiều ngược kim đồng hồ. Vì vậy tương tự
quay thước kẹp, khi liên tục xét điểm kế tiếp trên bao lồi ngoài, ta liên tục
duy trì cạnh gần nhất tương ứng trên bao lồi trong. Cách này có độ phức tạp
$O(n)$ cho một bộ dữ liệu.

Chứng minh rằng ``chọn điểm trên bao lồi ngoài tại một đỉnh không bao giờ tệ
hơn'' như sau. Nếu điểm được chọn nằm trên một cạnh của bao lồi ngoài:
\begin{enumerate}
  \item Nếu cạnh gần nhất của bao lồi trong song song với cạnh này, dời điểm
  đã chọn trên bao lồi ngoài đến một đầu mút sẽ không tệ hơn.
  \item Nếu cạnh gần nhất của bao lồi trong không song song với cạnh này, dời
  điểm đã chọn trên bao lồi ngoài đến một đầu mút sẽ tốt hơn.
\end{enumerate}
Suy ra chọn điểm tại một đỉnh của bao lồi ngoài không bao giờ tệ hơn.

\section{Smzzl with Tropical Taste}

Thể tích trà chanh trong bể thỏa phương trình vi phân
\[
  \begin{cases}
    V'(t)=(q-p)V(t),\\
    V(0)=1.
  \end{cases}
\]
Giải ra
\[
  V(t)=e^{(q-p)t}.
\]
Thực chất ta cần xét tích phân suy rộng
\[
  \int_{0}^{+\infty}V(t)\,dt
\]
có hội tụ hay không. Nguyên hàm là
\[
  \int V(t)\,dt=
  \begin{cases}
    \dfrac{e^{(q-p)t}}{q-p}, & q\ne p,\\
    t, & q=p.
  \end{cases}
\]
Vì vậy khi $q<p$ thì tích phân hội tụ; khi $q\ge p$ thì phân kỳ.

\section{Yiwen with Formula}

Gọi $\textit{cnt}_x$ là số dãy con có tổng bằng $x$. Khi đó đáp án là
\[
  \prod x^{\textit{cnt}_x}.
\]
Xét cách tính $\textit{cnt}_x$. Một cách đơn giản là đặt $f[i][j]$ là số cách
chọn từ $i$ số đầu để có tổng $j$, rồi làm DP với độ phức tạp $O(n^2)$. Để tối
ưu DP này, viết mỗi số $a_i$ thành đa thức $1+x^{a_i}$, sau đó dùng chia để
trị kết hợp NTT. Chú ý mô-đun của đa thức thực ra là
\[
  \varphi(998244353)=998244352,
\]
nên cần NTT modulo tùy ý. NTT ba mô-đun hoặc FFT tách bit không tối ưu có thể
bị quá thời gian vì hằng số quá lớn; bản mã đầu tiên của người ra đề thậm chí
chạy chậm hơn vét cạn. Có thể tham khảo phương pháp tối ưu do Mao Xiao nêu
trong luận văn đội tuyển quốc gia năm 2016.

Độ phức tạp là $O(n\log^2 n)$.

\section{Yiwen with Sqc}

Dễ thấy các chữ cái khác nhau không ảnh hưởng lẫn nhau, nên xét riêng từng chữ
cái. Không mất tính tổng quát, xét chữ cái $a$. Giả sử $a$ xuất hiện
\textit{cnt} lần trong chuỗi gốc. Gọi vị trí xuất hiện thứ $i$ của chữ cái $a$
trong chuỗi gốc là $a_i$, với $a_0=0$ và $a_{\textit{cnt}+1}=n+1$. Đặt
\[
  \textit{len}_i=a_{i+1}-a_i,
\]
tức khoảng cách giữa hai chữ cái $a$ kề nhau.

Trước hết cố định đầu trái tại $1$ và di chuyển đầu phải. Khi đó đáp án là
\[
  \sum_{k=1}^{\textit{cnt}}\textit{len}_k\,k^2,
\]
vì khi đầu phải di chuyển giữa hai vị trí $a_i$ kề nhau, số lượng chữ cái
không đổi; biểu thức trên chính là tổng theo từng đoạn, mỗi đoạn lấy độ dài
nhân với bình phương số chữ cái từ đầu đoạn đến đầu chuỗi.

Liên tục dịch đầu trái và gọi vị trí hiện tại là \textit{pos}. Khi
$\textit{pos}\le a_1$, đáp án đều bằng
\[
  \sum_{k=1}^{\textit{cnt}}\textit{len}_k\,k^2,
\]
nên đóng góp của đoạn này là
\[
  \textit{len}_0
  \sum_{k=1}^{\textit{cnt}}\textit{len}_k\,k^2.
\]
Tương tự, khi $\textit{pos}\in(a_1,a_2]$, đáp án là
\[
  \sum_{k=2}^{\textit{cnt}}\textit{len}_k\,(k-1)^2.
\]
Có thể thấy với $\textit{pos}_i\in(a_i,a_{i+1}]$ ta có
\[
  \sum_{k=i+1}^{\textit{cnt}}\textit{len}_k\,(k-i)^2.
\]
Đặt
\[
  \textit{ans}_i=
  \sum_{k=i+1}^{\textit{cnt}}\textit{len}_k\,(k-i)^2.
\]
Nếu không muốn đọc suy luận chặt chẽ, có thể tự lấy vài ví dụ rồi lấy sai phân
hai lần của \textit{ans} để nhìn ra cách làm.

Gọi $\textit{dif}_i$ là sai phân của hai giá trị \textit{ans} kề nhau. Khi đó
\[
\begin{aligned}
  \textit{dif}_i
  &= \textit{ans}_{i+1}-\textit{ans}_i \\
  &= \sum_{k=i+2}^{\textit{cnt}}\textit{len}_k\,(k-i-1)^2
     -
     \sum_{k=i+1}^{\textit{cnt}}\textit{len}_k\,(k-i)^2 \\
  &= \sum_{k=i+2}^{\textit{cnt}}\textit{len}_k\,(k-i-1)^2
     -
     \sum_{k=i+2}^{\textit{cnt}}\textit{len}_k\,(k-i)^2
     -\textit{len}_{i+1} \\
  &= -\sum_{k=i+2}^{\textit{cnt}}\textit{len}_k\,(2k-2i-1)
     -\textit{len}_{i+1}.
\end{aligned}
\]

Gọi $\textit{diff}_i$ là sai phân của hai giá trị \textit{dif} kề nhau. Khi đó
\[
\begin{aligned}
  \textit{diff}_i
  &= \textit{dif}_{i+1}-\textit{dif}_i \\
  &= -\sum_{k=i+3}^{\textit{cnt}}\textit{len}_k\,(2k-2(i+1)-1)
     -\textit{len}_{i+2}
     +\sum_{k=i+2}^{\textit{cnt}}\textit{len}_k\,(2k-2i-1)
     +\textit{len}_{i+1} \\
  &= -\sum_{k=i+3}^{\textit{cnt}}\textit{len}_k\,(2k-2(i+1)-1)
     +\sum_{k=i+3}^{\textit{cnt}}\textit{len}_k\,(2k-2i-1)
     +3\textit{len}_{i+2}
     +\textit{len}_{i+1}
     -\textit{len}_{i+2} \\
  &= \sum_{k=i+3}^{\textit{cnt}}2\textit{len}_k
     +3\textit{len}_{i+2}
     +\textit{len}_{i+1}
     -\textit{len}_{i+2}.
\end{aligned}
\]
Từ công thức trên có thể dễ dàng duy trì giá trị \textit{diff}, rồi duy trì
\textit{dif}, và cuối cùng duy trì \textit{ans}.

Cách làm trên là $O(n)$; cách $O(26n)$ có thể bị kẹt hằng số.

\end{document}
